\documentclass[12pt]{article}

\setlength{\textwidth}{16cm} \setlength{\textheight}{23cm}
\setlength{\oddsidemargin}{-0.0cm} \setlength{\topmargin}{-1cm}
\usepackage{amssymb, amsmath}
\usepackage{tabularx}
\usepackage{float}

\newcommand{\sdfrac}[2]{\mbox{\small$\displaystyle\frac{#1}{#2}$}}
\newcommand{\fdfrac}[2]{\mbox{\footnotesize$\displaystyle\frac{#1}{#2}$}}
\newcommand{\ltfrac}[2]{\mbox{\large$\frac{#1}{#2}$}}


%\usepackage[latin2]{inputenc}
\usepackage{chngcntr}
%\counterwithin{figure}{section}
\usepackage{latexsym}
\usepackage{amssymb,amsmath,amsopn}
\usepackage[dvips]{graphicx}   
\usepackage{xcolor,epsfig}         
\usepackage{wrapfig}		
\usepackage[a4paper, margin={2.0cm,2.5cm}]{geometry}
\usepackage{setspace}
\setstretch{1.05}
\usepackage{upgreek}
\usepackage{fancyhdr}
\usepackage{url}
%\usepackage{natbib}
\usepackage[sort, numbers]{natbib}
\usepackage[utf8]{inputenc}
\usepackage{titlesec}
\usepackage{amsthm}
%\usepackage{marvosym}
\usepackage{caption}
\usepackage{multirow}

\usepackage{ulem}


\newtheorem{remark}{Remark}[section]
\newtheorem{proposition}{Proposition}[section]
\newtheorem{corollary}{Corollary}[section]
\newtheorem{theorem}{Theorem}[section]
\newtheorem{lemma}{Lemma}[section]
\newtheorem{definition}{Definition}[section]

\numberwithin{equation}{section}

%\def\newtitlea#1{ \centerline{\Large{\bf #1}} }
%\def\newtitleb#1{ \centerline{\Large{\bf #1}} \bigskip}
\newcolumntype{Y}{>{\centering\arraybackslash}X}





\begin{document}
	
	\def\halmos{\rule{6pt}{6pt}}
	
	
	
	
	% words and formulas 
	\def\a{\alpha}
	\def\aux{auxiliary}
	%\def\al{ \frac{2}{M+m} }
	\def\alb{al\-ge\-braic}
	\def\alg{algorithm}
	\def\apl{application}
	\def\apr{approximat}
	\def\ass{assumption}
	\def\b{\beta}
	\def\bb{ {\cal B} }
	\def\bbd{ {\cal D} }
	\def\bbr{ {\cal R} }
	%\def\bbt{(BB^T)^{-1}}
	\def\bc{boundary condition}
	\def\bd{boundary}
	\def\bdy{boundary}
	\def\bdd{bounded}
	\def\bob{{\bf b}}
	\def\bp{{\bf p}}
	\def\bu{{\mathbf u}}
	\def\bh{{\bf h}}
	\def\bv{{\bf v}}
	\def\bww{{\bf w}}
	\def\bt{{\cal Z}^t}
	\def\bvp{boundary value problem}
	\def\ca{{\cal A}}
	\def\cf{{\cal F}}
	\def\ck{{\cal K}}
	\def\caf{C_{\a}}
	\def\cddot{\, : \,}
	\def\cg{conjugate gradient}
	\def\cgd{conjugate direction}
	\def\cgm{conju\-gate gra\-dient me\-thod}
	\def\cho{{\cal H}_0}
	\def\coe{coefficient}
	\def\comb{combination}
	\def\corr{corresponding}
	\def\cn{condition number}
	\def\cnd{condition}
	\def\cno{condition number}
	\def\cv{converg}
	%\def\cs{cross-section}
	%\def\ct{constant}
	\def\d{\delta}
	\def\dbc{Dirichlet boundary condition}
	\def\dev{{\rm dev\,}}
	%\def\df{difficulties}
	\def\dhom{D\cap {H}_0^{\bot}}
	\def\diff{differenti}
	\def\dim{di\-men\-sion}
	\def\din{xmethod}
	\def\dir{Dirichlet}
	\def\disc{discretiz}
	\def\disp{\displaystyle}
	\def\div{{\rm div}\, }
	%\def\dk{D^2}
	%\def\dkh{\tilde D^2}
	\def\dn{\delta_n}
	\def\dnbc{Dirichlet--Neumann boundary condition}
	\def\dr{Dirichlet}
	\def\ds{}	%\def\ds{\, d\sigma}
	\def\eh{\varepsilon_h}
	\def\ek{e_k}
	%\def\ell{elliptic}
	\def\eq{equation}
	\def\en{{\bf N}}
	\def\enp{{\bf N}^+}
	\def\ep{elasto-plastic}
	\def\er{{\bf R}}
	\def\ern{{\bf R}^N}
	\def\erp{{\bf R}^+}
	\def\esp{energy \ space}
	\def\est{estimate}
	\def\ev{eigenvalue}
	\def\fd{finite difference}
	\def\fe{finite element}
	\def\fdm{finite difference method}
	\def\fem{finite element method}
	%\def\g{\gamma}
	\def\ga{\gamma}
	\def\Ga{\Gamma}
	%\def\gb{\Gamma_\beta}
	%\def\gbb{_{\Gamma_\beta}}
	%\def\gd{Gateaux differentiable}
	\def\gdo{generalized \diff al \op}
	\def\gfem{gradient--finite element method}
	\def\gm{gradient method}
	\def\gr{gradient}
	\def\gv{\ov{g}}
	%\def\hc{H^1_{*}}
	\def\hd{H^1_D(\Om)}
	\def\hdiv{H^1_{div}}
	\def\hde{_{H^1_D}}
	\def\hdd{_{H^1_D}}
	\def\hddr{_{H^1_D(\Om)^r}}
	\def\hdg{hardening}
	\def\hdr{H^1_D(\Om)^r}
	\def\hed{H^1_{div}}
	\def\he{H^1_0(\Om)}
	\def\hee{_{H^1_0}}
	\def\hek{H^{1/2}(\Ga_N)}
	\def\hekp{H^{1/2}(\pa\Om)}
	\def\hi{H^1(\Om)}
	\def\hs{H^1_0(S)}
	\def\hh{H^2(\Om)\cap H^1_0(\Om)}
	\def\hhd{H^2(\Om)\cap H^1_D(\Om)}
	\def\hhn{ {H^4(\Om)\cap H^2_0(\Om)} }
	\def\hk{H^2(\Om)}
	\def\hn{h_n}
	\def\hnn{H^n_0(\Om)}
	\def\hnnr{H^n_0(\Om)^r}
	\def\ho{H_0}
	%\def\hom{H_0^{\bot}}
	\def\hq{H^1_D(\Om)}
	\def\hqr{H^1_D(\Om)^r}
	\def\hqq{_{H^1_D}}
	\def\hqqr{_{H^1_D(\Om)^r}}
	\def\hr{H^1_{rot}}
	\def\hsp{Hilbert space}
	\def\ip{inner product}
	\def\iter{iteration}
	\def\iv{^{-1}}
	\def\inv{investigat}
	\def\k{\kappa}
	\def\K{{\cal K}}
	\def\kaf{K_{\a}}
	\def\kkom{K_{2,\Om}}
	\def\kkgn{K_{2,\Ga_N}}
	\def\kkpao{K_{2,\pa\Om}}
	\def\kpom{K_{p,\Om}}
	\def\knom{K_{4,\Om}}
	\def\koetom{K_{5,\Om}}
	\def\kpeom{K_{{p_1},\Om}}
	\def\kpg{K_{p,\Ga}}
	\def\kpgn{K_{p,\Ga_N}}
	\def\kpkgn{K_{{p_2},\Ga_N}}
	\def\ks{^{(k-1)}}
	\def\kv{\ov K}        
	\def\kvh{K(\vh)}        
	\def\l{\lambda}
	\def\L{\Lambda}
	\def\la{\langle}
	\def\lap{\Delta}
	\def\las{linear algebraic system}
	\def\lk{L^2(\Om)}
	\def\lkb{L^2(B)}
	\def\lkgn{L^2(\Ga_N)}
	\def\lkk{_{L^2}}
	\def\li{linear}
	\def\lm{(-\Delta)^{-1}}
	%\def\m{^{-1}}
	\def\mt{method}
	%\def\n{^n}
	\def\nb{\nabla}
	\def\nbv{\overline\nabla}
	\def\nbhd{neighbourhood}
	\def\nequiv{\not\equiv}
	\def\nfem{Newton--finite element method}
	\def\nm{numerical}
	\def\no{\noindent}
	\def\nt{Newton}
	\def\ntm{Newton method}
	\def\ntsm{Newton's method}
	\def\om{\omega}
	\def\oml{\Omega_l}
	\def\Om{\Omega}
	\def\op{operator}
	\def\ov{\overline}
	%\def\p{{\cal P}}
	\def\pa{\partial}
	\def\pde{partial differential equation}
	\def\Po{Poisson}
	\def\pol{polynomial}
	\def\pos{posedness}
	\def\plg{ {\cal P}_{alg} }
	\def\pr{problem}
	\def\prad{ {\cal P}_{rad} }
	\def\prec{precondition}
	\def\ps{ {\cal P}_s }
	\def\pc{ {\cal P}_c }
	\def\pos{posedness}
	%\def\q{ \frac{M-m}{M+m} }
	\def\qn{ \frac{M_n-m_n}{M_n+m_n} }
	\def\ra{\rangle}
	\def\rad{radial}
	\def\real{realization}
	\def\resp{respectively}
	\def\rd{reaction-diffusion}
	\def\rn{{\bf R}^N}
	\def\rnn{{\bf R}^{N\times N}}
	\def\rnpr{{\bf R}^{(N+1)r}}
	\def\rnr{{\bf R}^{Nr}}
	\def\rot{{\rm rot}}
	\def\rr{{\bf R}^r}
	\def\rs{{\bf R}^s}
	\def\rss{{\bf R}^{s\times s}}
	\def\rth{\rightharpoonup}
	\def\Rtw{\Rightarrow}
	%\def\s{\sigma}
	\def\sbt{\subset}
	\def\sd{steepest descent}
	\def\sem{semilinear}
	\def\sl{solution}
	\def\slae{system of linear algebraic equations}
	\def\sn{S^{(n)}}
	\def\spd{strictly positive definite}
	\def\sq{sequence}
	\def\ssp{Sobolev space}
	\def\st{such that}
	\def\strf{stress function}
	\def\sy{system}
	%\def\t{^T}
	\def\Th{\Theta}
	\def\ti{\tilde}
	\def\tin{translation invarian}
	\def\tio{translation invariant with respect to $H_0$}
	\def\tn{\tau_n}
	\def\tor{torsion}
	\def\trf{transform}
	\def\trig{tri\-go\-no\-me\-tric}
	\def\tru{truncation}
	\def\ts{tangential stress}
	\def\tu{\tilde u}
	\def\ua{{\bf u}^\ast}
	\def\uas{{u}^\ast}
	\def\uh{{\bf u}_h}
	\def\uhh{u^{h}}
	\def\un{u_n}
	\def\unv{\ov u_n}
	\def\ux{x}
	\def\u#1{ u_{[#1]} }
	\def\ve{\varepsilon}
	\def\veu{\varepsilon({\bf u})}
	\def\vh{V_h}
	\def\vol{{\rm vol\,}}
	\def\vp{\varphi}
	\def\vr{\varrho}
	\def\ws{weak solution}
	\def\wti{\widetilde}
	\def\xc{x_{comp}}
	%\def\z{\zeta}
	\def\zn{z_n}
	\def\znv{\ov z_n}
	\def\zs{z_n^\ast}
	
	
	\def\sm{\smallskip}
	\def\bi{\bigskip}
	\def\me{\medskip}
	\def\ce{\centerline}
	
	
	
	\def\no{\noindent}
	\def\bi{\bigskip}
	\def\sm{\smallskip}
	\def\me{\medskip}
	
	\newcommand{\alertg}[1]{{\color{green}#1}}
	\newcommand{\alert}[1]{{\color{red}#1}}
	\newcommand{\alertb}[1]{{\color{blue}#1}}
	
	
	%xaxa
	
	\title{Remarks to a simplified nonlinear Stokes problem}
%	\author{J. Kar\'atson\footnote{Department
%			of Applied Analysis \& HUN-REN--ELTE Numerical Analysis and Large Networks Research Group, E\"otv\"os Lor\'and University; \quad
%			Department of Analysis and Operations Research, Budapest University of Technology and Economics, Hungary; \quad
%			kajkaat@caesar.elte.hu (corresponding author); \quad orcid: 0000-0003-1369-7743}, S. Sysala\footnote{Department of Applied Mathematics and Computational Sciences, Institute of Geonics of the Czech Academy of Sciences, Ostrava, Czechia; \quad stanislav.sysala@ugn.cas.cz; \quad orcid: 0000-0002-2704-4797}, M. B\'ere\v{s}\footnote{Department of Applied Mathematics and Computational Sciences, Institute of Geonics of the Czech Academy of Sciences, Czechia; Department of Applied Mathematics, VSB - Technical University of Ostrava, Czechia; michal.beres@ugn.cas.cz; \quad orcid: 0000-0001-8588-3268}}
	
	\maketitle

% \begin{abstract}
%    This paper is devoted to the application of a quasi-Newton/variable preconditioning (QN) approach to non-smooth problems, motivated by elasto-plastic models. Two versions are discussed: the first one is carried out via regularized approximations of the nonsmooth problem, and the second one gives an extension to non-smooth operators in order to be applied directly. Convergence analysis is presented for both variants.  Then these abstract methods are applied to elasto-plasticity where two different preconditioners are investigated. We also propose convenient iterative solvers for linearized systems appearing in each step of the QN method. The convergence results are illustrated on numerical examples in 3D inspired by real-life problems, and they demonstrate that the suggested QN methods are competitive with the standard Newton method. \alert{The corresponding Matlab's codes are available for download.}
% \end{abstract}
% 
%{\small 
%	
%\alert{	\noindent	
%	\textbf{Key words:}  nonlinear operator equation, non-smoothness, quasi-Newton/variable preconditioning, elasto-plasticity, large-scale systems.
%
%\me\noindent	
%	\textbf{MSC codes:} 35J60, 35Q74, 47J05, 65J15, 65N12, 74C05, 74S05
%	
%	\me\noindent	
%	\textbf{Declarations and statements:} The authors declare that they have no known competing financial interests or personal relationships that could have appeared to influence the work reported in this paper. We also declare that this paper is not submitted elsewhere.}
%}	
%	
%	\me
	
	
\section{Introduction}

In this report, we consider the quasi-Newton/variable preconditioning 
(QNVP) method for solving a simplified nonlinear Stokes problem. The simplification means that the pore pressure field is given \alert{and the constraint on divergence-free velocity field is not considered. Thus} the simplified problem has a same structure as nonlinear elasticity with the infinitesimal strain tensor.
	
First, we start with preliminaries to the QNVP method. Consider the nonlinear operator equation
\begin{equation}
	\label{opeq}
	F(u)= 0 \quad\mbox{in } H,
\end{equation}
where $H$ is a real Hilbert space with \ip \ $\la .,. \ra$ and induced norm $\|.\|$, $F:H\rightarrow H$ is a smooth operator and its derivative $F'$ has the following properties:
\begin{enumerate}
		\item[(i)] (Symmetry.) For any $u\in H$ the operator $F'(u)$ is self-adjoint.		
		\item[(ii)] (Ellipticity.) There exist constants $\Lambda\ge \lambda>0$  such that  
		\begin{equation}
		\label{f2}	
		\lambda \|v\|^2\leq\langle F'(u)v,v\rangle \leq \Lambda \|v\|^2\qquad (\forall u, v\in H).
		\end{equation}		
		\item[(iii)] (Lipschitz continuity.)	There exists a  constant  $L>0$ such that  
		$$
		\|F'(u)-F'(v)\|\leq L \, \|u-v\|  \qquad (\forall u,\, v\in H).
		$$
\end{enumerate}

The QNVP method enables us to solve the equation \eqref{opeq} by the following algorithm:
\begin{equation}
	\label{qniter}
	u^{(n+1)}:= u^{(n)} -\frac{2}{M_n+m_n} \, B_n\iv F(u^{(n)})\qquad (n\in\en), 
\end{equation}
where $B_n$ is a self-adjoint operator in $H$ and $m_n, M_n$ are positive constants such that $0<m\leq m_n\le M_n\leq M$ and
\begin{equation}
		\label{qnitsp}
		m_n \la B_n h,h \ra \le \la F'(u_n)h,h \ra \le M_n  \la B_n h,h \ra 
		\qquad  (n\in\en, \, h\in H).
\end{equation}
	

\section{The simplified nonlinear Stokes problem}
	\label{sec_EP}
	
The aim of this section is to specify and analyze the operators $F$ and $F'$ for the investigated problem. We consider a bounded domain $\Om\subset\er^3$ with a Lipschitz continuous boundary. The space $H$ is a finite-element approximation of the Sobolev space $[H^1_D(\Omega)]^3$ equipped with the following inner product and norm:
$$\langle u,v \rangle=\int_{\Omega}\varepsilon(u):\varepsilon(v),\quad \|v\|^2=\langle v,v \rangle.$$ The subscript $D$ means that Dirichlet-type boundary conditions are prescribed somewhere on $\Gamma_D$. The mapping $F$ is defined by
\begin{equation}
\langle F(u),v\rangle=\int_{\Omega}a(|\varepsilon(u_0+u)|)\varepsilon(u_0+u):\varepsilon(v)-\int_\Omega gv,\quad (u,v\in H),
\end{equation}  
where $u_0\in [H^1(\Omega)]^3$ and $g\in [L^2(\Omega)]^3$ are given functions representing nonhomogeneous Dirichlet boundary conditions and a pore pressure gradient, respectively. The function $a$ represents the Carreau law and is prescribed in the following form \cite{ABS22, JMS10}:
\begin{equation}
a(r)=\mu_\infty+(\mu_0-\mu_\infty)(1+\lambda r^2)^{\frac{p-2}{2}}, \quad r\geq0,
\label{a_function}
\end{equation}
where $p\in(1,2)$, $\lambda>0$, and $\mu_0>\mu_\infty>0$. \alert{[Comment: There is a small discrepancy in the literature. The formula \eqref{a_function} is in accordance with \cite{ABS22}. However, by \cite{JMS10}, $\lambda^2$ should be used instead of $\lambda$. this will complete us the choice of $\lambda$.]} Clearly, $a(0)=\mu_0$ and the values of $a$ belongs to the interval $(\mu_\infty,\mu_0]$. Since
\begin{equation}
	a'(r)=(\mu_0-\mu_\infty)(p-2)\lambda r(1+\lambda r^2)^{\frac{p-4}{2}}<0 \quad\forall r>0,
\end{equation}
the function $a$ is decreasing in $(0,+\infty)$. We also need to work with the function $\bar a(r):=a(r)+ra'(r)$. It can be found in the form
\begin{equation}
	\bar a(r):=\mu_\infty+(\mu_0-\mu_\infty)(1+(p-1)\lambda r^2)(1+\lambda r^2)^{\frac{p-4}{2}}, \quad r\geq0.
\end{equation}
We have: $\bar a(r)\leq a(r)$, $\bar a(0)=\mu_0$, and $\bar a(r)\in(\mu_\infty,\mu_0]$.
The graph comparing the functions $a$ and $\bar a$ for $\mu_0=1$, $\mu_\infty=1e-3$, $\lambda=1$, and $p=1.3$ is depicted in Figure \ref{fig_a_function}. 

\begin{figure}[H]
\centering
\includegraphics[width=0.6\textwidth]{a_function.png}
\caption{Comparison of the functions $a$ and $\bar a$ for $\mu_0=1$, $\mu_\infty=1e-3$, $\lambda=1$, and $p=1.3$.}
\label{fig_a_function}
\end{figure}

Using the functions $a$ and $\bar a$, one can subsequently derive the following formulas and estimates:
\begin{equation*}
	\langle F'(u)w,v\rangle=\int_{\Omega}\Big[a(|\varepsilon(u+u_0)|)\varepsilon(w):\varepsilon(v)+\frac{a'(|\varepsilon(u+u_0)|)}{|\varepsilon(u+u_0)|}(\varepsilon(u+u_0):\varepsilon(v))(\varepsilon(u+u_0):\varepsilon(w))\Big],
\end{equation*}
\begin{equation*}
	\langle F'(u)v,v\rangle=\int_{\Omega}\Big[a(|\varepsilon(u+u_0)|)|\varepsilon(v)|^2+|\varepsilon(u+u_0)|a'(|\varepsilon(u+u_0)|)\Big(\frac{\varepsilon(u+u_0)}{|\varepsilon(u+u_0)|}:\varepsilon(v)\Big)^2\Big],
\end{equation*}
\begin{equation}
	\int_{\Omega}\bar a(|\varepsilon(u+u_0)|)|\varepsilon(v)|^2\leq\langle F'(u)v,v\rangle\leq\int_{\Omega} a(|\varepsilon(u+u_0)|)|\varepsilon(v)|^2,
	\label{F'_bound}
\end{equation}
Comparing \eqref{f2} with \eqref{F'_bound}, we obtain
\begin{equation}
\lambda=\min_{r\geq0}\,\bar a(r)=\mu_\infty,\quad\Lambda=\max_{r\geq0}\,a(r)=\mu_0.
\end{equation}

	
%%%%%%%%%%%%%%%%%%%%%%%%%%%%%%%%%%%%%%%%% Section 3 %%%%%%%%%%%%%%%%%%%%%%%%%%%%%%%%%%%%%%%%%%%%%%%%%%%%%%%%%%%%%%
	
\section{Preconditioners $B_n$ for the magnetic problem}	
\label{sec_QN_plast}

In this section, we introduce variable and fixed preconditioners defining the algorithm \eqref{qniter}. It is convenient to initiate this algorithm with $u^{(0)}=0$.  

\subsection{Variable preconditioner}

The variable preconditioner $B_n$ for the investigated problem is in the form
\begin{equation}
\langle B_n w,v\rangle=\int_{\Omega}c(|\varepsilon(u^{(n)}+u_0)|)\varepsilon(w):\varepsilon(v),
\label{B_n}
\end{equation}
where 
\begin{equation}
	c(r)=(1-\gamma)a(r)+\gamma\bar a(r),\quad \gamma\in[0,1].
	\label{c1}
\end{equation} 
Clearly, $\bar a(r)\leq c(r)\leq a(r)$ for any $r\geq0$. Comparing \eqref{B_n} with \eqref{F'_bound}, one can easily derive that the estimate 
\eqref{qnitsp} holds with 
\begin{equation}
	\begin{array}{l}
m_n=\inf\limits_{x\in\Omega}\frac{\bar a(|\varepsilon((u^{(n)}+u_0)(x))|)}{c(|\varepsilon((u^{(n)}+u_0)(x))|)}=\frac{1}{(1-\gamma)\varrho_n+\gamma},\\[3mm]
M_n=\sup\limits_{x\in\Omega}\frac{a(|\varepsilon((u^{(n)}+u_0)(x))|)}{c(|\varepsilon((u^{(n)}+u_0)(x))|)}=\frac{\varrho_n}{(1-\gamma)\varrho_n+\gamma},
	\end{array}
	\label{Mm}
\end{equation}
where
\begin{equation}
	\varrho_n=\sup_{x\in\Omega}\frac{a(|\varepsilon((u^{(n)}+u_0)(x))|)}{\bar a(|\varepsilon((u^{(n)}+u_0)(x))|)}=\frac{M_n}{m_n}.
	\label{varrho}
\end{equation}
This ratio is a measure how accurately $B_n$ approximates $F'(\nabla u_n)$. It can be estimated as follows:
\begin{equation}
	\varrho_n\leq\sup_{r\geq0}\frac{a(r)}{\bar a(r)}\leq \sup_{r\geq0}\frac{1+\lambda r^2}{1+(p-1)\lambda r^2}=\lim_{r\rightarrow+\infty}\frac{1+\lambda r^2}{1+(p-1)\lambda r^2}=\frac{1}{p-1}.
	\label{varrho_bound}
\end{equation}
Here, the following implication was used:
$$b\geq c>\mu_\infty>0\qquad\Longrightarrow\qquad \frac bc\leq\frac{b-\mu_\infty}{c-\mu_\infty}.$$
Let us note that the second inequality in \eqref{varrho_bound} is sufficiently sharp if $\mu_\infty<<\mu_0$. Otherwise, this bound can be too pessimistic. Let us also note that $(M_n+m_n)/2=1$ if $\gamma=1/2$ in \eqref{c1}. This follows from \eqref{Mm}.

\subsection{Fixed preconditioner}

A fixed preconditioner $B:=B_n$ for the investigated problem can be found in the form
\begin{equation}
	\langle B_n w,v\rangle=\int_{\Omega}\mu_0\varepsilon(w):\varepsilon(v),
	\label{B_n2}
\end{equation}
where $\mu_0$ is a material parameter defining the function $a$. The corresponding parameters $m_n$ and $M_n$ can be estimated as follows:
\begin{equation}
	m_n=\frac{1}{\mu_0}\inf_{x\in\Omega}\bar a(|\varepsilon((u^{(n)}+u_0)(x))|),\quad M_n=\frac{1}{\mu_0}\sup_{x\in\Omega}a(|\varepsilon((u^{(n)}+u_0)(x))|).
	\label{Mm2}
\end{equation}
This preconditioner is convenient if we expect that the values $|\varepsilon((u^{(n)}+u_0)(x))|$ are sufficiently small for any $x\in\Omega$. In such a case, $a(|\varepsilon((u^{(n)}+u_0)(x))|)\approx\mu_0$ and $\bar a(|\varepsilon((u^{(n)}+u_0)(x))|)\approx\mu_0$.

\section{Numerical examples}

\subsection{Example 1: prismatic bar}
\label{subsec_ex1}

First, we consider a model problem with a prismatic bar occupying the domain $\Omega=(0,\rho)\times(0,\rho)\times(0,L)$, where $L$ is the length of the bar and $\rho<<L$. We assume that the bar is aligned with the axis $x_3$, i.e., $x_1,x_2\in(0,\rho)$ and $x_3\in(0,L)$. On the prism shell, the unknown vector function $u=(u_1,u_2,u_3)$ is zero in the normal direction \alert{(slip conditions)}. On the left face with $x_3=0$, $u$ is constant in the normal direction, i.e., $u_3(x_1,x_2,0)=u_{3,0}=const$. On the right face with $x_3=L$, the homogeneous Neumann boundary condition is prescribed.  The volume force $g$ is considered only in the direction $x_3$ and is constant, i.e., $g=(0,0,g_3)$, $g_3=const<0$. Finally, we assume a homogeneous material. 

\begin{remark}
Under the aforementioned assumptions, the solution satisfies $u=(0,0,u_3)$, where $u_3$ depends only on $x_3$ and satisfies the following differential equation:
$$a(|u_3'(x_3)|)u_3'(x_3)=-g_3(x_3-L),\quad u_3(0)=0.$$ 
This equation can be solved e.g. by the forward Euler method and used for validation of 3D numerical results presented below. If $p=2$ then $a(r)=\mu_0$ and
$$u_3(x_3)=-\frac{g_3}{2\mu_0}x_3^2+\frac{g_3}{\mu_0}Lx_3+u_{3,0}.$$
\alert{Let us also note that the unknown field $u=(0,0,u_3)$ is divergence-free if and only if $u_3$ is constant with respect to the coordinate $x_3$. This is possible only for $g_3=0$.}
\end{remark}


Now, we present numerical results for the following input data: $L=50$, $\rho=5$, $u_{3,0}=5$, $\mu_0=1$, $\mu_\infty=0.001$, $\lambda=10$, $p=1.1$, and several different values of $g_3$. Furthermore, we use P1 elements and a uniform mesh with 103,680 elements and 54,886 unknowns. As in \cite{KSB25}, we combine the algorithm \eqref{qniter} with the deflated and preconditioned conjugate gradient methods. The preconditioner is defined either by the incomplete Cholesky method or aggregation-based algebraic multigrid (AGMG) method.

In Figures \ref{fig_velocity} and \ref{fig_a_values}, the computed velocity field and the corresponding values of the function $a$ for $g_3=-0.003$ are depicted, respectively. We see that the velocity is decreasing along the length of the bar. The values of $a$ are close to $\mu_0=1$. The strongest nonlinearity is observed for $x_3=0$.
\begin{figure}[H]
	\centering
	\includegraphics[width=0.4\textwidth]{velocity.png}
	\caption{Computed velocity field for $g_3=-0.003$.}
	\label{fig_velocity}
\end{figure}
\begin{figure}[H]
	\centering
	\includegraphics[width=0.8\textwidth]{a_values.png}
	\caption{Values of the functions $a$ corresponding to the solution for $g_3=-0.003$.}
	\label{fig_a_values}
\end{figure}

Convergence of the Newton, quasi-Newton 1, and quasi-Newton 2 methods for $g_3=-0.003$ is compared in Figure \ref{fig_iteration}. The convergence is very fast for all the method. In case of the Newton method, superlinear convergence is observed.
\begin{figure}[H]
	\centering
	\includegraphics[width=0.8\textwidth]{iteration.png}
	\caption{Convergence of the Newton, quasi-Newton 1, and quasi-Newton 2 methods.}
	\label{fig_iteration}
\end{figure}

Convergence and runtime for the Newton (N), quasi-Newton1 (qN1), and quasi-Newton2 (qN2) methods are compared in Table \ref{table_ex1}. We consider several values of $g_3$, the incomplete Cholesky (IC) method, and the AGMG method. One can see that the convergence and runtime of the Newton method are not so much sensitive on $g_3$, unlike qN1 and qN2. The qN2 method is fastest for smaller values of $|g_3|$. However, for $|g_3|=0.007$ or larger, the convergence of qN2 is the slowest, as material response is strongly nonlinear. One can also see that IC is slightly faster than AGMG.
\alert{[S: My laptop has been used for computation.]}

\begin{table}[h!]
\caption{Convergence and runtime for the investigated solvers and Example 1. Values in the table are Newton's iterations/Cumulative CG iterations/Runtime [s].}
\begin{tabular}{|c||c|c|c||c|c|c|}
\hline
$g_3$ & IC-N & IC-qN1 & IC-qN2 & AGMG-N & AGMG-qN1 & AGMG-qN2 \\ \hline\hline
-3e-3 & 6/112/1.86 & 14/84/0.95 & 15/58/\textbf{0.61} & 6/51/2.12 & 14/35/0.99 & 15/44/0.85\\ \hline
-4e-3 & 6/111/1.84 & 17/90/1.17 & 22/70/\textbf{0.79} & 6/53/2.01 & 17/40/1.25 & 22/57/1.11\\ \hline
-5e-3 & 7/130/2.17 & 27/97/1.72 & 35/93/\textbf{1.12} & 7/63/2.44 & 25/50/1.85 & 36/86/1.68\\ \hline
-6e-3 & 7/133/2.17 & 41/111/2.69 & 74/167/\textbf{2.11} & 7/63/2.39 & 41/70/3.05 & 75/170/3.33\\ \hline
-7e-3 & 8/152/\textbf{2.53} & 73/143/5.17 & $>$200/-/5.53 & 8/76/2.84 & 72/108/5.63 & $>$200/-/6.81\\ \hline
\end{tabular}
\label{table_ex1}
\end{table}

\subsection{Example 2: prismatic bar with non-constant thickness}

Now, we consider a similar model problem as in Example 1, however, with a more complicated geometry. In particular, we assume a prismatic bar of the length $L$, which is aligned with the axis $x_3$, i.e., $x_3\in(0,L)$. Cross sections of the bar are squares with the size $\rho$ depending on the coordinate $x_3$ by the following formula: 
$\rho(x_3):=\rho_1+(\rho_2-\rho_1)x_3/L$, where $\rho_1, \rho_2>0$ are given values. It means that the computational domain can be defined by the following formula:
$$\Omega:=\big\{(x_1,x_2,x_3)\in\mathbb R^3\ |\; x_3\in(0,L),\; x_1,x_2\in\big(-\rho(x_3)/2,\rho(x_3)/2\big)\big\}.$$
On the prism shell, the unknown vector function $u=(u_1,u_2,u_3)$ satisfies $u_3=0$. \alert{[S: Please note that the slip condition would be in the form $u\cdot n=0$.]} On the left face with $x_3=0$, we assume that $u$ is constant in the normal direction, i.e., $u_3(x_1,x_2,0)=u_{3,0}=const$. On the right face with $x_3=L$, the homogeneous Neumann boundary condition is prescribed. The volume force $g$ is considered only in the direction $x_3$ and is constant, i.e., $g=(0,0,g_3)$, $g_3=const<0$. Next, we assume a homogeneous material and set the following input data: $L=50$, $\rho_1=6$, $\rho_2=5$, $u_{3,0}=5$, $\mu_0=1$, $\mu_\infty=0.001$, $\lambda=10$, $p=1.1$, and several different values of $g_3$. We use P1 elements and a uniform mesh with 86,400 elements and 50,986 unknowns. As in \cite{KSB25}, we combine the algorithm \eqref{qniter} with the deflated and preconditioned conjugate gradient methods. The preconditioner is defined either by the incomplete Cholesky method or AGMG method.

In Figures \ref{fig_velocity2} and \ref{fig_a_values2}, the computed velocity field and the corresponding values of the function $a$ for $g_3=-0.003$ are depicted, respectively. We see that the velocity is decreasing along the length of the bar. The values of $a$ are close to $\mu_0=1$. The strongest nonlinearity is observed for $x_3=0$.
\begin{figure}[H]
	\centering
	\includegraphics[width=0.4\textwidth]{velocity2.png}
	\caption{Computed velocity field for $g_3=-0.003$.}
	\label{fig_velocity2}
\end{figure}
\begin{figure}[H]
	\centering
	\includegraphics[width=0.4\textwidth]{a_values2.png}
	\caption{Values of the functions $a$ corresponding to the solution $g_3=-0.003$.}
	\label{fig_a_values2}
\end{figure}

Convergence of the Newton, quasi-Newton 1, and quasi-Newton 2 methods with respect to the AGMG preconditioner is compared in Figure \ref{fig_iteration2}. The convergence is very fast for all the method. In case of the Newton method, superlinear convergence is observed.
\begin{figure}[H]
	\centering
	\includegraphics[width=0.8\textwidth]{iteration2.png}
	\caption{Convergence of the Newton, quasi-Newton 1, and quasi-Newton 2 methods.}
	\label{fig_iteration2}
\end{figure}

Convergence and runtime for the Newton (N), quasi-Newton1 (qN1), and quasi-Newton2 (qN2) methods are compared in Table \ref{table_ex2}. We consider several values of $g_3$, the incomplete Cholesky (IC) method, and the AGMG method. One can see that the convergence and runtime of the Newton method are not so much sensitive on $g_3$, unlike qN1 and qN2. The qN2 method is fastest for smaller values of $|g_3|$. However, for $|g_3|=0.008$ or larger, the convergence of qN2 is the slowest, as material response is strongly nonlinear. One can also see that IC is faster than AGMG.
\alert{[S: My laptop has been used for computation.]}

\begin{table}[h!]
	\caption{Convergence and runtime for the investigated solvers and Example 2. Values in the table are Newton's iterations/Cumulative CG iterations/Runtime [s].}
	\begin{tabular}{|c||c|c|c||c|c|c|}
		\hline
		$g_3$ & IC-N & IC-qN1 & IC-qN2 & AGMG-N & AGMG-qN1 & AGMG-qN2 \\ \hline\hline
		-3e-3 & 5/86/1.28 & 12/72/0.76 & 14/56/\textbf{0.52} & 5/39/1.52 & 12/33/0.90 & 13/39/0.74\\ \hline
		-4e-3 & 6/100/1.57 & 16/77/0.93 & 17/59/\textbf{0.59} & 6/50/1.80 & 16/37/1.06 & 17/51/0.93\\ \hline
		-5e-3 & 6/101/1.56 & 19/83/1.15 & 23/72/\textbf{0.73} & 6/50/1.88 & 19/43/1.35 & 24/70/1.25\\ \hline
		-6e-3 & 7/119/1.85 & 25/90/1.56 & 36/97/\textbf{1.07} & 7/61/2.21 & 25/53/1.87 & 37/95/1.80\\ \hline
		-7e-3 & 7/120/1.87 & 38/102/2.32 & 68/162/\textbf{1.77} & 7/61/2.17 & 37/68/2.71 & 70/167/3.09\\ \hline
		-8e-3 & 8/139/\textbf{2.20} & 62/126/3.87 & 168/347/4.15 & 8/71/2.52 & 61/93/4.44 & 198/359/7.28\\ \hline
	\end{tabular}
	\label{table_ex2}
\end{table}

%%%%%%%%%%%%%%%%%%%%%%%%%
\section{Remarks to generalized Poiseuille flow problem}

%\alert{Janos, you mentioned that Example 1 from Section \ref{subsec_ex1} could be similar or even the same as the Poiseuille flow problem if we replace the slip boundary conditions on the prism shell with the non-slip ones. First of all, I have tried to implement Example 1 with the non-slip boundary conditions. It led to a very strange velocity field, which practically vanished for $x_3>u_{3,0}$. In  have found that this is caused by the absence of the divergence-free constraint on the velocity field. After that, I have studied the Poiseuille flow problem and found that it is defined without the divergence-free constraint, but only in 2D.}

In this auxiliary section, we introduce a generalized Poiseuille flow problem. The generalization means that the nonlinear function $a$ defined by \eqref{a_function} is considered. We show that it suffices to solve an auxiliary 2D problem to find its solution. Next, we eliminate the unknown pressure and the div-constraint from the original problem. Finally, we derive a problem, where $\varepsilon(u)$ is used instead of $\nabla u$.

Let $\Omega$ be a cylindrical domain defined by 
$$\Omega:=\{(x_1,x_2,x_3)\in\mathbb R^3\ |\; (x_1,x_2)\in\widetilde\Omega,\;x_3\in(0,L)\},$$
where $\widetilde\Omega$ is a Lipschitz domain in 2D. Let $p_1, p_2>0$, $p_1<p_2$, be given constant pressures prescribed at the bottom and top faces of $\Omega$, respectively. Then the generalized (steady-state) Poiseuille flow problem reads: \textit{find} $u\colon\Omega\rightarrow\mathbb R^3$ and $p\colon\Omega\rightarrow\mathbb R$ such that
\begin{equation}
\left\{\begin{array}{l}
	-\mathrm{div}[a(|\nabla u|)\nabla u]+\nabla p=0\;\;\mbox{in }\Omega,\\
	\mathrm{div}\,u=0\;\;\mbox{in }\Omega,\\
	u=0 \;\;\mbox{on }\{(x_1,x_2,x_3)\in\mathbb R^3\ |\; (x_1,x_2)\in\partial\widetilde\Omega,\;x_3\in(0,L)\},\\
	-a(|\nabla u|)[\nabla u] n+pn=(0,0,-p_1)^\top \;\;\mbox{on }\{(x_1,x_2,x_3)\in\mathbb R^3\ |\; (x_1,x_2)\in\widetilde\Omega,\;x_3=0\},\\
	-a(|\nabla u|)[\nabla u] n+pn=(0,0,p_2)^\top \;\;\mbox{on }\{(x_1,x_2,x_3)\in\mathbb R^3\ |\; (x_1,x_2)\in\widetilde\Omega,\;x_3=L\}.
\end{array}
\right.
\label{P1}
\end{equation}
The solution of \eqref{P1} satisfies: 
\begin{equation}
p(x_1,x_2,x_3)=p_1+(p_2-p_1)x_3/L,
\label{p}
\end{equation}
\begin{equation}
u=(u_1,u_2,u_3),\quad u_1=u_2=0, \quad u_3=u_3(x_1,x_2),
\label{u}
\end{equation}
\begin{equation}
	\left\{\begin{array}{c}
		-\mathrm{div}[a(|\nabla u_3|)\nabla u_3]=(p_1-p_2)/L=:g_3\;\;\mbox{in }\widetilde\Omega,\\
		u_3=0 \;\;\mbox{on }\partial\widetilde\Omega.
	\end{array}
	\right.
\label{u_3}
\end{equation}

\begin{proof}
We have: $-\nabla p=(0,0,g_3)$ and
$$\nabla u=\left(
\begin{array}{ccc}
0 & 0 & 0\\
0 & 0 & 0\\
\partial u_3/\partial x_1 & \partial u_3/\partial x_2 & 0
\end{array}
\right),\quad |\nabla u|=|\nabla u_3|,\quad \mathrm{div}\,u=0.$$
The outward normal $n$ is equal to $(0,0,-1)$ on the bottom face and $(0,0,1)$ on the top face. Therefore, \eqref{p}--\eqref{u_3} imply \eqref{P1}$_1$--\eqref{P1}$_5$. 
\end{proof}

Now, we eliminate $p$ and $\mathrm{div}\,u=0$ from the problem \eqref{P1}. After that, we arrive at the following reduced problem:
\textit{find} $u\colon\Omega\rightarrow\mathbb R^3$ such that
\begin{equation}
	\left\{\begin{array}{l}
		-\mathrm{div}[a(|\nabla u|)\nabla u]=(0,0,g_3)^\top\;\;\mbox{in }\Omega,\\
		u=0 \;\;\mbox{on }\{(x_1,x_2,x_3)\in\mathbb R^3\ |\; (x_1,x_2)\in\partial\widetilde\Omega,\;x_3\in(0,L)\},\\
		a(|\nabla u|)[\nabla u] n=0 \;\;\mbox{on }\{(x_1,x_2,x_3)\in\mathbb R^3\ |\; (x_1,x_2)\in\widetilde\Omega,\;x_3=0\},\\
		a(|\nabla u|)[\nabla u] n=0 \;\;\mbox{on }\{(x_1,x_2,x_3)\in\mathbb R^3\ |\; (x_1,x_2)\in\widetilde\Omega,\;x_3=L\}.
	\end{array}
	\right.
	\label{P2}
\end{equation}
Clearly, the solution $u$ satisfies \eqref{u} and \eqref{u_3}. Furthermore, if \eqref{u} and \eqref{u_3} hold then 
$$\varepsilon(u)=\frac{1}{2}(\nabla u+(\nabla u)^\top)=\left(
\begin{array}{ccc}
	0 & 0 & \frac12\frac{\partial u_3}{\partial x_1}\\
	0 & 0 & \frac12\frac{\partial u_3}{\partial x_2}\\
	\frac12\frac{\partial u_3}{\partial x_1} & \frac12\frac{\partial u_3}{\partial x_2} & 0
\end{array}
\right),\quad |\varepsilon(u)|=|\nabla u|/\sqrt{2}.$$
Hence, one can rewrite the problem \eqref{P2} as follows: \textit{given} $u_3\colon\Omega\rightarrow\mathbb R$ satisfying \eqref{u_3}, \textit{find} $u\colon\Omega\rightarrow\mathbb R^3$ such that
\begin{equation}
	\left\{\begin{array}{l}
		-\mathrm{div}[a(\sqrt2|\varepsilon(u)|)\varepsilon(u)]=(0,0,g_3)^\top\;\;\mbox{in }\Omega,\\
		u=0 \;\;\mbox{on }\{(x_1,x_2,x_3)\in\mathbb R^3\ |\; (x_1,x_2)\in\partial\widetilde\Omega,\;x_3\in(0,L)\},\\
		a(\sqrt2|\varepsilon(u)|)[\varepsilon(u)] n=-\frac{1}{2}a(|\nabla u_3|)(\partial u_3/\partial x_1, \partial u_3/\partial x_2, 0)^\top \;\;\mbox{on bottom face},\\
		a(\sqrt2|\varepsilon(u)|)[\varepsilon(u)] n=\;\;\:\frac{1}{2}a(|\nabla u_3|)(\partial u_3/\partial x_1, \partial u_3/\partial x_2, 0)^\top \;\;\mbox{on top face}.
	\end{array}
	\right.
	\label{P3}
\end{equation}

\alert{Janos, if we wish to relate Example 1 from Section \ref{subsec_ex1} with the Poiseuille flow problem, then it is necessary to solve the problem \eqref{u_3} at first and then change the boundary conditions at the bottom and top faces.}


%\alert{[S: I am not sure whether $\varepsilon(u)$ or $\nabla u$.]}

%%%%%%%%%%%%%%%%%%%%%%%%%
\section{Planned work -- recommendation}

\alert{The simplified nonlinear Stokes problem studied in this paper is rather nonlinear elasticity. It is not too useful for solving flow problems. I recommend to complete the theory of QNVP method for the saddle-point problems at first and then to start to investigate nonlinear incompressible flow problems. Alternatively, one can try to combine Uzawa's algorithm or other splitting procedures with QNVP.}


	
\begin{thebibliography}{99}
	
\bibitem{ABS22}		
Anguiano, M., Bonnivard, M., \& Suárez-Grau, F. J. (2022). Carreau law for non-Newtonian fluid flow through a thin porous media. The Quarterly Journal of Mechanics and Applied Mathematics, 75(1), 1-27.

\bibitem{JMS10}	
Janela, J., Moura, A., \& Sequeira, A. (2010). A 3D non-Newtonian fluid?structure interaction model for blood flow in arteries. Journal of Computational and applied Mathematics, 234(9), 2783-2791.
			
\bibitem{KF03}		
Karátson, J., \& Faragó, I. (2003). Variable preconditioning via quasi-Newton methods for nonlinear problems in Hilbert space. SIAM journal on numerical analysis, 41(4), 1242-1262.

\bibitem{KSB25}
Karatson, J., Sysala, S., \& Béreš, M. (2025). Quasi-Newton iterative solution approaches for nonsmooth elliptic operators with applications to elasto-plasticity. Computers \& Mathematics with Applications, 178, 61-80.
					

\end{thebibliography}
		
		
\end{document}
	
	
	
	